\documentclass[10pt,a4paper]{article}

\usepackage[utf8]{inputenc}
\usepackage[T1]{fontenc}

\usepackage{amsmath,mathtools,amsfonts,amssymb}
\usepackage{graphicx}
\usepackage[spanish]{babel}
\usepackage{lastpage,tabto,xcolor}
\usepackage[a4paper, top=0.8in, bottom=1in, left=0.7in, right=0.7in]{geometry}
\usepackage{fancyhdr}
\usepackage{comment}
\usepackage{multicol}
\usepackage{titling} % Para ajustar los espacios del título
\usepackage{tasks} % Para listas en columnas
\usepackage{tikz}
\usepackage{pgfplots}
\pgfplotsset{width=7cm,compat=1.15}
\usepackage{caption}

% Fecha de versión automática según fecha de compilación
\newcommand{\verdate}{\the\year.\ifnum\month<10 0\fi\the\month.\ifnum\day<10 0\fi\the\day}

% Ajusta la separación entre el encabezado y el contenido
\setlength{\headsep}{10pt}

% Ajuste del título:
\setlength{\droptitle}{-0.5in} % Mueve el título ligeramente hacia arriba
\pretitle{\begin{center}\vspace{-1em}\normalfont\Large}
	\posttitle{\par\end{center}\vspace{-4em}}
\preauthor{\begin{center}\vspace{-0.5em}}
	\postauthor{\par\end{center}}
\predate{}\postdate{}

\pagestyle{fancy}
\lhead{\footnotesize Análisis de Señales y Sistemas  - Año \the\year{} Ver. \verdate}
\rhead{\footnotesize Trabajo Práctico N$^\text{o}$4 -- Parte A: Serie de Fourier de tiempo continuo.}
\rfoot{\footnotesize Página \thepage\ de \pageref{LastPage}}
\renewcommand{\headrulewidth}{0.4pt}
\renewcommand{\footrulewidth}{0.4pt}

% Título optimizado con espacios reducidos
\title{%
	{\normalsize Departamento de Ingeniería en Tecnologías Electrónicas, UTN FRM}\\[0.1cm]
	{\large Análisis de Señales y Sistemas}\\[0.15cm]
	{\normalsize Trabajo Práctico N$^{\text o}$4 -- Parte A: Serie de Fourier de tiempo continuo}%
}
\author{}
\date{}


\begin{document}
	\maketitle
	\thispagestyle{fancy}

	% Entorno de dos columnas
	\begin{multicols}{2}

		\label{sec:enun}
		% ***********************************************
			\textbf{Producto interno y ortogonalidad}
			\begin{enumerate}
			% ***********************************************
			% 				Ejercicio 1
			% ***********************************************
			\item Sea $\omega_0 = 2\pi/T$. Calcular por cálculo directo, sobre un período de longitud $T$, el producto interno
			\[
				\langle e^{jm\omega_0 t},\, e^{jn\omega_0 t}\rangle
			\]
			y concluir sobre la ortogonalidad del conjunto $\{e^{jk\omega_0 t}\}_{k\in\mathbb{Z}}$. Particularizar el cálculo para $m=3$, $n=5$ y $T = 2\pi$.

			% ***********************************************
			% 				Ejercicio 2
			% ***********************************************
			\item Calcular la norma cuadrática sobre un período de longitud $T = 2\pi/\omega_0$ de las siguientes señales:
			\begin{enumerate}
				\item $e^{jk\omega_0 t}$, con $k\in\mathbb{Z}$.
				\item $\cos(k\omega_0 t)$, con $k\geq 1$.
				\item $\sin(k\omega_0 t)$, con $k\geq 1$.
			\end{enumerate}

			% ***********************************************
			\textbf{Serie de Fourier exponencial}
			% ***********************************************
			% 				Ejercicio 3
			% ***********************************************
			\item Considerar el pulso rectangular periódico
			\[
				x(t) =
				\begin{cases}
					1, & -\tfrac{\tau}{2} < t < \tfrac{\tau}{2}, \\
					0, & \tfrac{\tau}{2} < |t| < \tfrac{T}{2},
				\end{cases}
				\qquad T = 3\tau.
			\]
			\begin{enumerate}
				\item Calcular los coeficientes $c_k$ de la serie exponencial e identificar la envolvente.
				\item Graficar el espectro de líneas (módulo y fase) para $0 \leq k \leq 4$ y completar el rango $-4 \leq k < 0$ aprovechando la paridad de $x(t)$.
			\end{enumerate}

			% ***********************************************
			% 				Ejercicio 4
			% ***********************************************
			\item Considerar la señal triangular periódica definida como $x(t) = 1 - |t|$ para $-1 < t < 1$, con $T = 2$.
			\begin{enumerate}
				\item Calcular los coeficientes $c_k$ de la serie exponencial.
				\item Graficar el espectro de líneas (módulo y fase) para $0 \leq k \leq 4$ y completar el rango $-4 \leq k < 0$ aprovechando la paridad de $x(t)$.
			\end{enumerate}

			% ***********************************************
			% 				Ejercicio 5
			% ***********************************************
			\item Considerar el tren de impulsos periódico de período $T_0$,
			\[
				\delta^{T_0}(t) = \sum_{k=-\infty}^{\infty} \delta(t - kT_0).
			\]
			Calcular los coeficientes $c_k$ de la serie exponencial de Fourier y graficar el espectro de líneas.

			% ***********************************************
			\textbf{Propiedades de la serie de Fourier}
			% ***********************************************
			% 				Ejercicio 6
			% ***********************************************
			\item Sea $x(t)$ el pulso rectangular periódico del Ej. 3, con coeficientes $c_k$ ya calculados. Aplicando la propiedad de desplazamiento temporal, obtener los coeficientes $\tilde{c}_k$ de la señal desplazada $\tilde{x}(t) = x(t - \tau/2)$ sin reintegrar. Graficar el nuevo espectro de fases para $0 \leq k \leq 4$ y completar el rango $-4 \leq k < 0$ aprovechando que $\tilde{x}(t)$ es real. Comparar con el espectro original y comentar las diferencias.

			% ***********************************************
			% 				Ejercicio 7
			% ***********************************************
			\item Para la señal triangular del Ej. 4:
			\begin{enumerate}
				\item Calcular la potencia media de $x(t)$ por integración directa sobre un período.
				\item Estimar la potencia media usando la identidad de Parseval truncada en $|k| \leq 5$, es decir $\sum_{k=-5}^{5} |c_k|^2$.
				\item Indicar la fracción de la potencia total contenida en los primeros cinco armónicos.
			\end{enumerate}

			% ***********************************************
			\textbf{Serie trigonométrica}
			% ***********************************************
			% 				Ejercicio 8
			% ***********************************************
			\item Sea la onda cuadrada periódica de amplitud $\pm 1$ y período $T = 4$:
			\[
				x(t) =
				\begin{cases}
					-1, & -2 < t < 0, \\
					\phantom{-}1, & \phantom{-}0 < t < 2.
				\end{cases}
			\]
			\begin{enumerate}
				\item Identificar por simetría qué coeficientes trigonométricos resultan nulos sin necesidad de calcularlos.
				\item Calcular $b_k$.
				\item Obtener los coeficientes $c_k$ de la serie exponencial aplicando las relaciones entre coeficientes trigonométricos y exponenciales.
			\end{enumerate}

			% ***********************************************
			% 				Ejercicio 9
			% ***********************************************
			\item Partiendo de los coeficientes $c_k$ del Ej. 3 (pulso rectangular) y del Ej. 4 (triangular), obtener los coeficientes trigonométricos $a_k$ y $b_k$ aplicando las relaciones inversas. Comentar cómo se manifiesta la paridad de cada señal en los coeficientes obtenidos.

			% ***********************************************
			\textbf{Serie de Fourier y sistemas LIT}
			% ***********************************************
			% 				Ejercicio 10
			% ***********************************************
			\item Sea el sistema LIT con respuesta al impulso $h(t) = e^{-t}u(t)$.
			\begin{enumerate}
				\item Calcular la respuesta en frecuencia $H(j\omega)$ a partir de la definición.
				\item Determinar la salida en régimen sinusoidal permanente ante la entrada $x(t) = \cos(2t)$, indicando módulo y fase de $H(j\omega)$ en la frecuencia correspondiente.
				\item Tomar como entrada periódica la señal triangular del Ej. 4 ($\omega_0 = \pi$) y, usando los coeficientes $c_k$ ya calculados, obtener los coeficientes $d_k$ de la salida aplicando $d_k = H(jk\omega_0)\,c_k$. Graficar los espectros de líneas (módulo) de entrada y salida para $0 \leq k \leq 4$ y completar el rango $-4 \leq k < 0$ aprovechando que $h(t)$ es real (módulo par en $k$).
			\end{enumerate}

			% ***********************************************
			% 				Ejercicio 11
			% ***********************************************
			\item Considerar el sistema LIT con respuesta en frecuencia
			\[
				H(j\omega) =
				\begin{cases}
					1, & |\omega| < 3, \\
					0, & |\omega| \geq 3.
				\end{cases}
			\]
			Graficar $H(j\omega)$ y, utilizando métodos del dominio de la frecuencia, encontrar la salida del sistema ante la entrada $x(t) = \cos(2t) + \sin(4t)$.

			% ***********************************************
			% 				Ejercicio 12
			% ***********************************************
			\item Sea un filtro pasabanda ideal con respuesta en frecuencia
			\[
				H(j\omega) =
				\begin{cases}
					1, & \omega_1 < |\omega| < \omega_2, \\
					0, & \text{en otro caso.}
				\end{cases}
			\]
			La entrada al filtro es el pulso rectangular periódico del Ej. 3, con $\omega_0 = 2\pi/(3\tau)$. Determinar la salida $y(t)$ y graficar su espectro de líneas (módulo) para los dos casos siguientes:
			\begin{enumerate}
				\item $\omega_1 = 0{,}5\,\omega_0$ y $\omega_2 = 2{,}5\,\omega_0$.
				\item $\omega_1 = 3{,}5\,\omega_0$ y $\omega_2 = 6{,}5\,\omega_0$.
			\end{enumerate}
			Comparar con el espectro de la entrada e indicar qué armónicos sobreviven en cada caso. Aprovechar la paridad de $x(t)$ para reducir el cálculo a $k \geq 0$ y reflejar al rango negativo.

		\end{enumerate}

	\end{multicols}

\end{document}
